\scanpage{13}
2-connexe des faisceaux en groupes, munis d'un isom. de \uline{noyaux} avec $K$, \uline{et} d'un sous-faisceau en groupes dans la classe indiquée par $\unsure{\varphi|\phi}$.

Si $c(K,\varphi)\in H^2(X,\underline{\unsure{\operatorname{Nor}|\operatorname{Na}}}(\varphi))'$, alors $c(K)\in H^2(X,\unsure{I|Z}(K))'$ !

[Principe de \unsure{réduction} de \unsure{Schur|Schreier|schéma}]

Autre principe de réduction. Supposons que nous ayons un hom \inserted{``surjectif''} de noyaux $K\longrightarrow K'$, et que nous ayons donné un isom $K'\simeq\widetilde{F}'$, où $F'$ est un faisceau en groupes.

\deleted{Alors il existe une} Considérons la catégorie locale des faisceaux en groupes $F$, muni d'un \unsure{hom} $F-F'$, et d'un isom $\widetilde{F}\xrightarrow{\sim}K$, \unsure{compatibles|compatible} avec $K\longrightarrow K'$ et $K'\simeq\widetilde{F}'$. C'est une catégorie 2-connexe, et correspond à un noyau $K''$ tel que l'on ait une suite exacte
\diagram{\begin{tikzcd}[ampersand replacement=\&]
e \arrow[r,dash] \& K'' \arrow[r,dash] \& I(K) \arrow[r] \& I(K') \arrow[r,dash] \& e
\end{tikzcd}}
\deleted{Alors} Considérons la classe $c'\in H^2(\unsure{k|X},K'')$. Son image dans \deleted{$\unsure{I(K}$} $H^2(k,I(K))$ est définie et est $c(K)$, donc si $c'\in H^2(X,K'')'$, alors $c\in H^2(X,I(K))'$.
% archive sheet number: 12

